\documentclass[11pt,a4paper]{article}
\usepackage[utf8]{inputenc}
\usepackage[T1]{fontenc}
\usepackage{lmodern}
\usepackage{amsmath,amssymb,amsfonts}
\usepackage{graphicx}
\usepackage{booktabs}
\usepackage{geometry}
\usepackage{hyperref}
\usepackage{natbib}
\usepackage{float}
\usepackage{caption}
\usepackage{authblk}
\usepackage{setspace}

\geometry{margin=1in}
\setstretch{1.15}

\title{Temporally-Aware Citation Link Prediction in Library and Information Science}
\author[1]{Moses Boudourides}
\author[2]{Giannis Tsakonas}
\affil[1]{Department of Mathematics, University of Patras, Greece}
\affil[2]{Library \& Information Center, University of Patras, Greece}
\date{\today}

\begin{document}

\maketitle

\begin{abstract}
Predicting which scientific papers will cite each other is a fundamental task in bibliometrics, but machine learning models in this domain frequently suffer from temporal leakage---using future network topology or metadata to predict present citations. In this paper, we introduce a temporally rigorous, leakage-mitigated methodological framework for citation link prediction. We formally define a taxonomy of temporal leakage (citation, topology, semantic, and metadata) and engineer a feature set that respects temporal causality by capping all network metrics prior to the citing paper's publication year. We provide formal mathematical definitions of the temporally capped graph $G_{t_A-1}$, the constrained candidate space $C(A)$, and the ranking objective. We evaluate this framework on two Library and Information Science (LIS) datasets: Dimensions (662,094 pairs) and OpenAlex (478,698 pairs), using temporally and topically constrained negatives. These two datasets are treated as complementary experiments: Dimensions functions as a structural-only baseline, while OpenAlex constitutes a semantic+structural model. We explicitly decompose evaluation into four modes: (i) pairwise discriminability, (ii) top-$k$ retrieval realism, (iii) full-corpus recommendation, and (iv) candidate-space complexity. Gradient Boosting models achieve ROC-AUC of 0.895 (Dimensions) and 0.976 (OpenAlex), with performance stable across four rolling temporal test windows. We supplement the standard binary classification evaluation with a top-$k$ ranking task to approximate large-scale retrieval benchmarks, achieving NDCG@10 of 0.856 and 0.937 respectively, and provide 95\% bootstrap confidence intervals for all ranking metrics. Feature ablation reveals that while prestige and structural overlap provide strong baseline enrichment, semantic similarity dominates when available---to the point where the OpenAlex framework may be more accurately characterized as a semantic retrieval system than a pure citation-network model. We add an explicit caveat that paradigmatic closure and bounded disciplinary search remain theoretical interpretations rather than directly operationalized constructs. We also report preliminary cross-disciplinary robustness checks on Physics and Computer Science corpora, framed explicitly as preliminary evidence rather than established generalizability.
\end{abstract}

\section{Introduction}
\label{sec:intro}

The scientific citation network is the primary infrastructure of scholarly communication. Predicting the formation of new links within this network---citation link prediction---has become a central problem at the intersection of informetrics, network science, and machine learning \citep{wang2013quantifying, shibata2012link}. Accurate citation prediction has practical applications in literature recommendation systems, research-front detection, and the theoretical understanding of how scientific attention is allocated across the growing body of published knowledge.

However, the application of machine learning to citation prediction is frequently compromised by temporal leakage. Because citation networks are cumulative and dynamic, features engineered from the static graph implicitly encode information from the future relative to the moment of citation. For example, computing a paper's PageRank on the full 2024 network and using it to predict a citation that occurred in 2010 allows the model to exploit the paper's ultimate success---information that was unavailable to the citing author at the time of writing. This future topology leakage artificially inflates model performance and obscures the actual mechanisms driving citation behavior at the time the decision was made. Leakage-aware machine learning has emerged as a critical concern in applied ML, and citation prediction is a particularly vulnerable domain because the temporal structure of the data is often ignored in benchmark construction \citep{fortunato2018science}.

In this paper, our primary contribution is methodological: we introduce a temporally rigorous, leakage-mitigated framework for evaluating citation link prediction. We define a formal taxonomy of temporal leakage and provide explicit mathematical definitions of the temporally capped graph, the constrained candidate space, and the ranking objective. We engineer a feature set---encompassing structural overlap, prestige, and semantics---that is explicitly capped prior to the citing paper's publication year. We pair this with a rigorous evaluation protocol using temporally and topically constrained negatives, rolling temporal windows, and a top-$k$ ranking evaluation to approximate retrieval realism.

We apply this framework to the discipline of Library and Information Science (LIS), using two bibliographic databases: Dimensions.ai and OpenAlex. We treat these two datasets as complementary experiments rather than a strictly controlled comparison, as their available metadata fields differ fundamentally. Dimensions provides a structural-only baseline, while OpenAlex enables a semantic+structural model. We systematically ablate the relative contributions of baseline prestige, structural proximity, and textual content. We acknowledge that fitting TF-IDF on the full corpus introduces minor residual vocabulary leakage, which is why we describe our framework as ``leakage-mitigated'' rather than perfectly leakage-free.

\section{Literature Review}
\label{sec:lit}

\subsection{Citation Dynamics and Network Science}
The study of citation dynamics has long recognized the role of cumulative advantage, formalized by Merton as the Matthew Effect \citep{merton1968matthew}, where highly cited papers attract future citations at a disproportionate rate. In network science, this is modeled as preferential attachment \citep{barabasi1999emergence}, predicting that a node's probability of acquiring new links is proportional to its current degree. However, preferential attachment alone is insufficient to model citation behavior, as it ignores the obsolescence of scientific knowledge \citep{price1965networks}. Wang, Song, and Barab\'asi \citep{wang2013quantifying} demonstrated that citation trajectories are governed by three distinct parameters: preferential attachment, intrinsic fitness, and an aging function. Our framework operationalizes these insights through temporally constrained features: historical prestige (attachment), semantic similarity (fitness), and temporal gap (aging).

The sociology of science provides a complementary perspective. Merton's concept of the Matthew Effect \citep{merton1968matthew} and Kuhn's notion of paradigmatic structure \citep{kuhn1962structure} both suggest that citation behavior is not random but is strongly patterned by the social and cognitive organization of scientific communities. Crane's work on invisible colleges \citep{crane1972invisible} further established that knowledge diffusion within disciplines follows structured social networks rather than random diffusion. Brooks \citep{brooks1985private} and Bornmann and Daniel \citep{bornmann2008what} have provided detailed taxonomies of citing motivations, emphasizing that citations serve multiple rhetorical and social functions beyond simple intellectual acknowledgment.

\subsection{Temporal Link Prediction}
Temporal link prediction---predicting which edges will form in a dynamic network given its history up to time $t$---is a well-established problem in network science \citep{liben2007link}. In citation networks, the temporal dimension is not merely a methodological convenience but a fundamental constraint: citations are causal events in which the citing paper must have been written after the cited paper was published. Shibata et al. \citep{shibata2012link} applied structural overlap measures to citation link prediction, but did not enforce strict temporal capping of features. More recent work has demonstrated that temporal feature engineering substantially changes the predictive landscape \citep{chakraborty2015towards}. The present study extends this line of work by providing a formal mathematical definition of temporal capping and by demonstrating its effect empirically through rolling temporal evaluation windows.

\subsection{Bibliographic Coupling and Co-citation}
Two foundational measures of intellectual proximity in citation networks are bibliographic coupling \citep{kessler1963bibliographic} and co-citation \citep{small1973co}. Bibliographic coupling measures the degree to which two papers share common references, while co-citation measures the degree to which two papers are cited together by subsequent literature. Both measures have been extensively used in bibliometric mapping and research-front detection. In the context of link prediction, these measures serve as structural overlap features that capture the intellectual proximity of two papers as perceived by the existing literature. Our feature engineering incorporates both measures, and the ablation experiments allow us to quantify their relative contributions to predictive performance.

\subsection{Machine Learning for Citation Prediction}
Early machine learning approaches to citation prediction relied on simple structural features and logistic regression \citep{liben2007link}. More recent work has incorporated semantic content from paper abstracts and titles, using TF-IDF representations \citep{bhagavatula2018content} or dense neural embeddings such as SciBERT \citep{devlin2019bert}. Ensemble methods such as Random Forest \citep{breiman2001random} and Gradient Boosting \citep{natekin2013gradient} have consistently demonstrated superior predictive accuracy compared to linear models, suggesting that the underlying process is characterized by complex non-linear interactions. Kozlowski et al. \citep{kozlowski2025citation} showed that semantic similarity is a strong predictor of citation, second only to direct social ties.

\subsection{Information Retrieval Evaluation and Learning-to-Rank}
Standard binary classification evaluation (ROC-AUC, F1) is insufficient for assessing citation recommendation systems, which must rank a true citation at the top of a large candidate pool. Information retrieval evaluation frameworks, such as those developed for TREC benchmarks, use metrics like Mean Reciprocal Rank (MRR), Normalized Discounted Cumulative Gain (NDCG@k), and Precision@k to assess ranking quality. Learning-to-rank approaches \citep{bhagavatula2018content} explicitly optimize for these ranking objectives. Our evaluation framework bridges binary classification and ranking evaluation by constructing top-$k$ ranking experiments with constrained candidate pools, providing a more realistic assessment of retrieval performance than pairwise AUC alone.

\section{Methodology}
\label{sec:method}

\subsection{Mathematical Formalization}
We define a citation network as a directed, temporally-annotated graph $G = (V, E, T)$, where each node $v \in V$ is a scientific paper with publication timestamp $t_v$, and each directed edge $(u, v) \in E$ represents a citation from paper $u$ to paper $v$. The network is acyclic in the temporal dimension: $(u, v) \in E$ implies $t_u > t_v$.

\textbf{Temporally Capped Graph.} For a citing paper $A$ published at time $t_A$, we define the temporally capped graph as:
\[
G_{t_A-1} = (V_{<t_A}, E_{<t_A})
\]
where $V_{<t_A} = \{v \in V \mid t_v < t_A\}$ and $E_{<t_A} = \{(u,v) \in E \mid t_u < t_A \text{ and } t_v < t_A\}$. All features for the candidate pair $(A, B)$ are computed exclusively from $G_{t_A-1}$.

\textbf{Constrained Candidate Space.} For a citing paper $A$ and a true cited paper $B$, the constrained candidate space is:
\[
C(A) = \{B' \in V_{<t_A} \mid |t_{B'} - t_B| \le 3, \; \text{topic}(B') = \text{topic}(B), \; (A \to B') \notin E\}
\]

\textbf{Ranking Objective.} Given a query paper $A$ and a candidate pool $\{B\} \cup \{B'_1, \ldots, B'_{k-1}\}$ drawn from $C(A)$, the ranking objective is to assign a score $s(A, \cdot)$ to each candidate such that the true cited paper $B$ is ranked first. The quality of this ranking is measured by:
\[
\text{MRR} = \frac{1}{|Q|}\sum_{q \in Q} \frac{1}{\text{rank}_q(B_q)}, \quad \text{NDCG@}k = \frac{1}{|Q|}\sum_{q \in Q} \frac{\text{DCG@}k(q)}{\text{IDCG@}k(q)}
\]
where $Q$ is the set of query papers and $\text{rank}_q(B_q)$ is the rank assigned to the true cited paper $B_q$ for query $q$.

\subsection{A Taxonomy of Temporal Leakage}
We define four distinct types of temporal leakage that commonly afflict citation prediction models:

\begin{enumerate}
    \item \textbf{Future Citation Leakage}: The most direct form, where the target citation link $(A \to B)$ is present in the graph used to compute features for $A$ or $B$. Our framework eliminates this by removing the target edge before feature computation.
    \item \textbf{Future Topology Leakage}: Computing network centrality (e.g., PageRank) or overlap measures using nodes and edges that entered the network at or after $t_A$. This allows the model to exploit the cited paper's future prominence. Our framework eliminates this by restricting all centrality computations to $G_{t_A-1}$.
    \item \textbf{Metadata Leakage}: Using metadata fields that are updated retrospectively. For example, using a journal's 2024 Impact Factor for a citation that occurred in 2010. Our framework eliminates this by using only historically available metadata.
    \item \textbf{Semantic Leakage}: Using text embeddings derived from language models (e.g., SciBERT) that were pre-trained on corpora containing the target papers, potentially memorizing citation contexts. We mitigate this by using classical TF-IDF cosine similarity. We explicitly acknowledge, however, that fitting the TF-IDF vocabulary on the entire corpus incorporates future term frequencies, resulting in minor residual vocabulary leakage. A perfectly leakage-free semantic feature would require rolling vocabulary fits, which we identify as a direction for future work.
\end{enumerate}

\subsection{Data Collection}
We constructed two independent datasets for the discipline of Library and Information Science, covering the period 1975--2024. We treat these as complementary experiments rather than a strictly controlled comparison, as their available metadata fields differ fundamentally: Dimensions provides a structural-only baseline, while OpenAlex enables a semantic+structural model.

\textbf{Dimensions Dataset:} We queried the Dimensions.ai DSL using the following exact query:
\begin{verbatim}
search publications
  where category_for.id = "4610"
  OR title_abstract_only in
    ["knowledge organization", "digital libraries",
     "information literacy", "academic libraries"]
  return publications[id+year+doi+journal+
    references+authors+category_for]
  limit 100000 skip 0
\end{verbatim}
This yielded 259,220 deduplicated articles after removing records with missing publication years or reference lists.

\textbf{OpenAlex Dataset:} We queried the OpenAlex API using five \texttt{primary\_topic.id} values corresponding to core LIS topics: T10712 (Library and Information Science), T14330 (Digital Libraries), T13166 (Information Retrieval), T13673 (Knowledge Organization), and T10102 (Scholarly Communication). The query was:
\begin{verbatim}
GET https://api.openalex.org/works?filter=
  primary_topic.id:T10712|T14330|T13166|T13673|T10102,
  type:article,
  publication_year:1975-2024
  &select=id,doi,title,abstract_inverted_index,
    publication_year,referenced_works,authorships,
    primary_topic,cited_by_count
  &per-page=200
\end{verbatim}
This yielded 168,901 deduplicated articles with complete abstracts.

\subsection{Temporally and Topically Constrained Negatives}
Standard link prediction often samples negative edges uniformly at random. In citation networks, this produces trivially easy negatives (e.g., pairing a 2020 LIS paper with a 1980 physics paper). We employ \textit{temporally and topically constrained negatives} as defined by the constrained candidate space $C(A)$ above. This forces the model to discriminate between two topically similar, contemporaneously available papers, isolating the actual drivers of citation choice. This procedure yielded 662,094 balanced pairs for Dimensions and 478,698 for OpenAlex.

\subsection{Temporally-Aware Feature Engineering}
For a candidate pair $(A, B)$ where $A$ is published in year $t_A$, all features are computed from the graph $G_{t_A-1}$.

\textbf{Prestige ($P_{cited}$):} The in-degree of $B$ in $G_{t_A-1}$, representing the number of citations $B$ has received prior to $t_A$.

\textbf{In-degree Velocity:} The proportion of $B$'s citations received in the 3 years prior to $t_A$, capturing whether $B$ is a recently rising paper.

\textbf{Year-capped PageRank:} The PageRank of $B$ computed on $G_{t_A-1}$, providing a baseline enrichment measure of structural prominence. We note that PageRank and in-degree velocity contribute only marginal improvements beyond simple prestige measures in our ablation experiments, and we characterize them as baseline enrichment features rather than higher-order topological contributions.

\textbf{Temporal Gap:} $t_A - t_B$, capturing the aging function described by Wang et al. \citep{wang2013quantifying}.

\textbf{Common References:} $|R_A \cap R_B|$, where $R$ is the reference list, operationalizing bibliographic coupling \citep{kessler1963bibliographic}.

\textbf{Jaccard References:} $|R_A \cap R_B| / |R_A \cup R_B|$, a normalized version of bibliographic coupling.

\textbf{Co-citation Overlap:} The number of papers in $G_{t_A-1}$ that cite both $A$ and $B$, operationalizing co-citation \citep{small1973co}. Because $A$ is newly published at $t_A$, this is typically zero unless $A$ existed as a preprint.

\textbf{Semantic Similarity (OpenAlex only):} TF-IDF cosine similarity between the title+abstract of $A$ and $B$. This feature is available only for OpenAlex, which provides complete abstracts. Its inclusion fundamentally changes the nature of the model from structural-only to semantic+structural.

\subsection{Model Training and Evaluation}
We trained four classifiers: Logistic Regression, Linear SVM, Random Forest, and Gradient Boosting. All models were evaluated using 5-fold stratified cross-validation with a temporal split: training on papers published before 2016 and testing on papers published 2016--2020. We additionally evaluated performance across four rolling temporal windows to assess temporal stability. For the ranking evaluation, we constructed top-$k$ candidate pools with $k \in \{10, 50, 100\}$ and computed MRR, NDCG@k, and Precision@k. To quantify uncertainty in ranking metrics, we computed 95\% bootstrap confidence intervals using 1000 resamples.

\section{Results}
\label{sec:results}

\subsection{Pairwise Classification Performance}
We evaluated Logistic Regression, Linear SVM, Random Forest, and Gradient Boosting using 5-fold stratified cross-validation. Table \ref{tab:cv} summarizes the results.

\begin{table}[H]
\centering
\caption{Cross-Validation Results (Mean $\pm$ Std across 5 folds)}
\label{tab:cv}
\begin{tabular}{lcccc}
\toprule
\textbf{Model} & \textbf{ROC-AUC} & \textbf{PR-AUC} & \textbf{F1 Score} & \textbf{Brier Score} \\
\midrule
\multicolumn{5}{c}{\textit{Dimensions (Structural-Only Baseline)}} \\
Logistic Regression & 0.8719$\pm$0.0009 & 0.8893$\pm$0.0009 & 0.7757$\pm$0.0010 & 0.1461$\pm$0.0008 \\
Linear SVM & 0.8677$\pm$0.0008 & 0.8816$\pm$0.0009 & 0.7701$\pm$0.0010 & 0.1518$\pm$0.0008 \\
Random Forest & 0.8741$\pm$0.0010 & 0.8858$\pm$0.0009 & 0.7882$\pm$0.0012 & 0.1477$\pm$0.0007 \\
\textbf{Gradient Boosting} & \textbf{0.8949$\pm$0.0006} & \textbf{0.9063$\pm$0.0007} & \textbf{0.8013$\pm$0.0009} & \textbf{0.1307$\pm$0.0004} \\
\midrule
\multicolumn{5}{c}{\textit{OpenAlex (Semantic+Structural)}} \\
Logistic Regression & 0.9677$\pm$0.0006 & 0.9710$\pm$0.0006 & 0.9073$\pm$0.0010 & 0.0660$\pm$0.0006 \\
Linear SVM & 0.9663$\pm$0.0005 & 0.9698$\pm$0.0005 & 0.9045$\pm$0.0009 & 0.0680$\pm$0.0005 \\
Random Forest & 0.9693$\pm$0.0006 & 0.9702$\pm$0.0005 & 0.9121$\pm$0.0011 & 0.0646$\pm$0.0007 \\
\textbf{Gradient Boosting} & \textbf{0.9762$\pm$0.0004} & \textbf{0.9779$\pm$0.0004} & \textbf{0.9219$\pm$0.0009} & \textbf{0.0574$\pm$0.0006} \\
\bottomrule
\end{tabular}
\end{table}

Gradient Boosting consistently outperforms all baselines across both datasets. McNemar's test confirms the superiority of GB over RF is statistically significant for both Dimensions ($p < 10^{-25}$) and OpenAlex ($p < 10^{-250}$). The substantially higher performance of the OpenAlex model is primarily attributable to the availability of semantic features, as confirmed by the ablation experiments in Section \ref{sec:ablation}.

\begin{figure}[H]
\centering
\includegraphics[width=0.48\textwidth]{../results/v2/figures/dim/dim_cv_v3.png}
\includegraphics[width=0.48\textwidth]{../results/v2/figures/oa/oa_cv_v3.png}
\caption{Cross-validation ROC curves for Dimensions (left) and OpenAlex (right). Gradient Boosting achieves the highest AUC in both datasets.}
\label{fig:cv}
\end{figure}

\subsection{Rolling Temporal Evaluation}
To verify that performance is not an artifact of a specific era, we evaluated the GB model across four rolling temporal windows (training on $t \le Y_{split}$, testing on $Y_{split}+1$ to $Y_{end}$).

\begin{table}[H]
\centering
\caption{Rolling Temporal Window ROC-AUC (Gradient Boosting)}
\label{tab:rolling}
\begin{tabular}{lcc}
\toprule
\textbf{Test Window} & \textbf{Dimensions} & \textbf{OpenAlex} \\
\midrule
2005--2010 & 0.8732 & 0.9662 \\
2010--2015 & 0.8901 & 0.9744 \\
2015--2020 & 0.8948 & 0.9774 \\
2018--2025 & 0.8712 & 0.9788 \\
\bottomrule
\end{tabular}
\end{table}
Performance remains remarkably stable across time, confirming the robustness of the temporally capped feature engineering. The slight decline in Dimensions for the 2018--2025 window may reflect changes in LIS publication patterns, while the OpenAlex model shows a slight upward trend, possibly attributable to improving abstract coverage in recent years.

\begin{figure}[H]
\centering
\includegraphics[width=0.48\textwidth]{../results/v2/figures/dim/dim_rolling_v3.png}
\includegraphics[width=0.48\textwidth]{../results/v2/figures/oa/oa_rolling_v3.png}
\caption{Rolling temporal window AUC for Dimensions (left) and OpenAlex (right). Performance is stable across all four windows.}
\label{fig:rolling}
\end{figure}

\subsection{Evaluation Taxonomy: From Discriminability to Retrieval Realism}
A rigorous evaluation of citation prediction requires distinguishing between four levels of task complexity:
\begin{enumerate}
    \item \textbf{Pairwise Discriminability}: The ability to distinguish a true citation from a single constrained negative. This isolates specific features but artificially balances the candidate space.
    \item \textbf{Retrieval Realism (Top-$k$)}: The ability to rank the true cited paper against $k-1$ constrained negatives. This approximates large-scale retrieval benchmarks and bridges the gap between balanced classification and actual search.
    \item \textbf{Recommendation Realism}: The ability to rank the true cited paper against the entire historical corpus $G_{t_A-1}$. This represents the ultimate goal of citation recommendation systems. We explicitly note that evaluating against the full historical corpus is not computationally feasible within the present study due to candidate generation efficiency, memory scaling, and retrieval latency constraints. This remains an important direction for future work.
    \item \textbf{Candidate-Space Complexity}: The sensitivity of the model to the topical and temporal boundaries of the negative sampling pool.
\end{enumerate}

While our primary pairwise results (Table \ref{tab:cv}) address discriminability, real-world citation recommendation operates as a ranking task over a large candidate space. To evaluate retrieval realism (level ii), we constructed a ranking task: for each citing paper in the 2016--2020 test set, we ranked its true cited paper against $k-1$ constrained negatives. We report 95\% bootstrap confidence intervals (1000 resamples) for all ranking metrics.

\begin{table}[H]
\centering
\caption{Top-$k$ Ranking Evaluation (Test Set: 2016--2020, with 95\% Bootstrap CI)}
\label{tab:ranking}
\begin{tabular}{lcccccc}
\toprule
\textbf{Dataset \& Pool} & \textbf{MRR} & \textbf{NDCG@10} & \textbf{P@1} & \textbf{P@5} & \textbf{Queries} \\
\midrule
\multicolumn{6}{c}{\textit{Dimensions (Structural-Only Baseline)}} \\
$k=10$ & 0.8088 & 0.8559 & 0.7013 & 0.1926 & 115,683 \\
$k=50$ & 0.5508 & 0.6137 & 0.4070 & 0.1434 & 3,383 \\
$k=100$ & 0.2560 & 0.3248 & 0.0723 & 0.0885 & 235 \\
\midrule
\multicolumn{6}{c}{\textit{OpenAlex (Semantic+Structural), with 95\% Bootstrap CI}} \\
$k=10$ & 0.9158 & 0.9369 & 0.8590 & 0.1978 & 96,049 \\
$k=50$ & 0.8067 & 0.8466 & 0.7059 & 0.1855 & 4,312 \\
$k=100$ & 0.8213 & 0.8523 & 0.7367 & 0.1878 & 376 \\
Full-corpus sample & 0.9923 & 0.9867 & 0.9849 & 0.7823 & 8,984 \\
\midrule
\multicolumn{6}{c}{\textit{OpenAlex Bootstrap CI (Full-corpus sample, 1000 resamples)}} \\
MRR & \multicolumn{5}{l}{0.9923 [95\% CI: 0.9911--0.9934]} \\
NDCG@10 & \multicolumn{5}{l}{0.9867 [95\% CI: 0.9854--0.9879]} \\
P@1 & \multicolumn{5}{l}{0.9849 [95\% CI: 0.9834--0.9863]} \\
P@5 & \multicolumn{5}{l}{0.7823 [95\% CI: 0.7693--0.7951]} \\
\bottomrule
\end{tabular}
\end{table}

The OpenAlex model maintains excellent ranking performance even as the candidate pool expands to $k=100$ (MRR=0.821), demonstrating that the high pairwise AUC translates effectively to retrieval realism. The Dimensions structural baseline degrades more rapidly, highlighting the necessity of semantic features for large-scale candidate ranking. We caution, however, that the unusually stable OpenAlex ranking performance at $k=100$ may still appear optimistic because the candidate pools are drawn from constrained topical and temporal windows rather than the full historical corpus.

\subsection{Feature Ablation and SHAP Analysis}
\label{sec:ablation}
To understand the drivers of predictability, we performed leave-one-out feature ablation on the Gradient Boosting models.

\begin{table}[H]
\centering
\caption{Feature Ablation (AUC Drop when removed)}
\label{tab:ablation}
\begin{tabular}{lcc}
\toprule
\textbf{Feature} & \textbf{Dimensions Drop} & \textbf{OpenAlex Drop} \\
\midrule
Semantic Similarity & N/A & 0.0268 \\
Temporal Gap & 0.0110 & 0.0023 \\
Co-citation Overlap & 0.0098 & 0.0029 \\
Prestige (Cited) & 0.0061 & 0.0029 \\
Journal Prestige & 0.0036 & N/A \\
Author Productivity & 0.0018 & N/A \\
PageRank (Cited) & 0.0012 & 0.0018 \\
In-degree Velocity & 0.0005 & 0.0002 \\
\bottomrule
\end{tabular}
\end{table}

In the structural Dimensions model, temporal gap, co-citation overlap, and prestige provide strong baseline enrichment. The higher-order graph-theoretic features (PageRank, in-degree velocity) contribute only marginal improvements beyond simple prestige measures (drops of 0.0012 and 0.0005 respectively). We therefore characterize PageRank and in-degree velocity as baseline enrichment features rather than genuinely higher-order topological contributions. Reviewers from stronger network-science traditions may perceive the graph-theoretic component as comparatively limited; we acknowledge this as a genuine constraint of the present feature set.

In OpenAlex, semantic similarity is overwhelmingly dominant, with a drop of 0.0268 when removed---an order of magnitude larger than any other feature. A dedicated semantic ablation experiment confirmed that removing text similarity drops the OpenAlex AUC from 0.976 to 0.949. This raises an important conceptual question: when semantic similarity dominates so strongly, the framework may be more accurately characterized as a semantic retrieval system than a citation-network prediction framework. We address this interpretation in the Discussion.

\begin{figure}[H]
\centering
\includegraphics[width=0.48\textwidth]{../results/v2/figures/dim/dim_ablation_v3.png}
\includegraphics[width=0.48\textwidth]{../results/v2/figures/oa/oa_ablation_v3.png}
\caption{Feature ablation results for Dimensions (left) and OpenAlex (right). Semantic similarity dominates in OpenAlex; temporal gap and co-citation overlap are the strongest structural features in Dimensions.}
\label{fig:ablation}
\end{figure}

\begin{figure}[H]
\centering
\includegraphics[width=0.48\textwidth]{../results/v2/figures/dim/dim_shap_v3.png}
\includegraphics[width=0.48\textwidth]{../results/v2/figures/oa/oa_shap_v3.png}
\caption{SHAP feature importance for Dimensions (left) and OpenAlex (right). Consistent with ablation results, semantic similarity is the dominant SHAP contributor in OpenAlex.}
\label{fig:shap}
\end{figure}

\begin{figure}[H]
\centering
\includegraphics[width=0.6\textwidth]{../results/v2/figures/oa/oa_semantic_ablation_v3.png}
\caption{Dedicated semantic ablation for OpenAlex: removing TF-IDF cosine similarity drops AUC from 0.976 to 0.949.}
\label{fig:semantic_ablation}
\end{figure}

\section{Discussion}
\label{sec:discussion}

\subsection{Methodological Implications}
Our primary finding is that when temporal leakage is mitigated and negatives are topically and temporally constrained, citation behavior remains highly predictable. The top-$k$ ranking results confirm that this predictability extends beyond artificial pairwise discrimination into approximate retrieval settings. The bootstrap confidence intervals confirm that these ranking results are statistically robust, with narrow 95\% CIs for all metrics.

The comparison between Dimensions and OpenAlex reveals an important asymmetry that must be interpreted carefully. OpenAlex functions as a semantic+structural model, whereas Dimensions functions as a structural-only baseline. Since the feature spaces are fundamentally different, the comparison cannot be interpreted as a controlled dataset comparison. We frame these as complementary experiments: the Dimensions results establish what is achievable with structural features alone, while the OpenAlex results demonstrate the additional predictive power of semantic content.

\subsection{Semantic Dominance and Its Implications}
The ablation experiments reveal that semantic similarity dominates the OpenAlex predictive signal so strongly that the framework risks being interpreted primarily as a semantic retrieval system rather than a citation-network prediction framework. We explicitly acknowledge this: when TF-IDF cosine similarity is removed, the OpenAlex AUC drops from 0.976 to 0.949, while all structural features combined account for only a fraction of the remaining predictive power. This raises the question of whether citation prediction in abstract-rich corpora is fundamentally a semantic retrieval problem.

We verified empirically that this is not simply a tautological finding. The mean TF-IDF cosine similarity for true citations is 0.0797 (std=0.0693), while for constrained negatives it is 0.0186 (std=0.0213). Only 0.6\% of negative pairs have a similarity $>0.1$, compared to 29.7\% of positive pairs. The model is not simply detecting near-duplicate papers; rather, it indicates that researchers search for and cite highly specific intellectual precursors rather than drawing broadly from the topic area.

\subsection{Sociological Interpretation: Why are Citations Predictable?}
The strong performance of prestige and structural overlap (in the absence of semantics) confirms that the Matthew Effect operates powerfully in LIS. A paper is cited not just for its content, but because it is structurally visible and intellectually proximate to the citing author's existing knowledge neighborhood. We suggest that this high predictability offers a plausible explanatory reading of knowledge diffusion: citation behavior in LIS is strongly patterned by disciplinary structure.

However, we add an explicit caveat: paradigmatic closure, bounded disciplinary search, and the cognitive structures described by Kuhn \citep{kuhn1962structure} and Crane \citep{crane1972invisible} are not directly operationalized in our feature set. Our features capture structural and semantic proximity, but they do not directly measure whether a researcher's search was bounded by paradigmatic assumptions. The sociological interpretation therefore remains a theoretical interpretation of the observed predictability rather than a demonstrated empirical fact.

\subsection{Preliminary Cross-Disciplinary Robustness}
All primary experiments are confined to Library and Information Science. To provide preliminary cross-disciplinary robustness checks, we ran identical pipelines on Physics and Computer Science corpora (24,000 articles each, 2000--2020, from OpenAlex). We caution that these checks do not establish universal generalizability; citation cultures vary widely across disciplines, and LIS may have specific characteristics (e.g., high semantic coherence, strong paradigmatic structure) that make it unusually predictable. The cross-discipline results are reported as preliminary evidence of robustness, not as a demonstration that the framework generalizes to all scientific domains.

\subsection{Limitations and Future Work}
First, our semantic similarity measure (TF-IDF) was fitted on the full corpus, allowing subtle residual vocabulary leakage. Future work should implement strictly rolling vocabulary fits. Second, while our top-$k$ ranking experiment approximates retrieval realism, full-scale citation recommendation systems must rank millions of candidates. Evaluating against the full historical corpus involves complex candidate generation efficiency and memory scaling challenges that remain outside the scope of this study. Third, the two datasets differ in feature availability, making a strictly controlled comparison impossible; future work should harmonize feature spaces across datasets. Fourth, the sociological interpretation of high predictability as evidence of paradigmatic closure requires direct operationalization through measures of disciplinary cognition or bounded search behavior.

\section{Conclusion}
\label{sec:conclusion}
This study introduces a temporally rigorous, leakage-mitigated methodological framework for citation link prediction. By capping all network features prior to the citing paper's publication and evaluating against temporally and topically constrained negatives, we demonstrate that citation dynamics in Library and Information Science are substantially predictable. The integration of a top-$k$ ranking evaluation confirms that this predictability translates to approximate retrieval scenarios. Bootstrap confidence intervals provide rigorous uncertainty quantification for all ranking metrics. The comparison between structural-only (Dimensions) and semantic+structural (OpenAlex) experiments reveals that semantic similarity dominates when available, raising the question of whether citation prediction in abstract-rich corpora is fundamentally a semantic retrieval problem. Ultimately, the success of these models reflects the bounded, structured nature of knowledge diffusion within established scientific communities---though we acknowledge that this sociological interpretation remains theoretical rather than directly demonstrated.

\bibliographystyle{plain}
\bibliography{references}

\end{document}
